\documentclass[11pt]{article}
\usepackage[margin=1.25in]{geometry}
\usepackage{microtype}
\pagestyle{empty}

\begin{document}

\noindent
Peter Cotton\\
\texttt{peter.cotton@gsmc.ai}\\[2em]

\noindent
The Editors\\
\emph{Economics Letters}\\[2em]

\noindent
Dear Editors,

\medskip

Please consider the enclosed manuscript, ``Winning the large language capability battle and losing the production economy,'' for publication in \emph{Economics Letters}.

The note makes a single observation and develops its consequence. When compute is itinerant, free to move between prediction tasks and to run any available algorithm, competitive equilibrium hands each subcontractible task to whichever model creates the most value per token of compute; capability enters only through that quotient. We exhibit a closed-form equilibrium in which a model strictly more accurate than every rival on every task nonetheless has production compute share tending to zero as the division of labor refines.

The measurement result may be of independent interest: for a sole frontier supplier at competitive scale, average PnL per token equals the shadow price of compute divided by the elasticity of task value with respect to effective compute, so the popular quotient ranks productivity exactly when that elasticity declines with crowding. Concavity alone is not sufficient, and we give a smooth counterexample within the stated assumptions.

The manuscript is not under consideration elsewhere. I have no conflicts of interest to declare. No data were used.

\medskip

\noindent
Sincerely,\\[1em]
Peter Cotton

\end{document}
